%% kapa: exam=1 round=33 year=2022 period=2 form=A subject=accounting subject_ko=회계학 seq=25 number=65
%% source: Q-Net 감정평가사 기출문제 (cst00309) artlSeq=5209997 file=raw/1008/5209997/2022년도 제33회 감정평가사 1차시험 2교시(감정평가 관계법규, 회계학).pdf
%% conversion: kordoc 4.24.51 -> tools/gen1.py (auto); review=proofread-sonnet; flags=none
\begin{question}{65}
(주)감평은 20×1년 1월 1일 사용목적으로 ￦5,000에 건물(내용연수 5년, 잔존가치 ￦0, 정액법 감가상각)을 취득하고 재평가모형을 적용하고 있다. 건물을 사용함에 따라 재평가잉여금 중 일부를 이익잉여금으로 대체하고, 건물 처분 시 재평가잉여금 잔액을 모두 이익잉여금으로 대체하는 정책을 채택하고 있다. 20×2년 말 건물에 대한 공정가치는 ￦6,000이다. (주)감평이 20×5년 1월 1일 동 건물을 처분할 때, 재평가잉여금 중 이익잉여금으로 대체되는 금액은?
\begin{choices}
  \item ￦0
  \item ￦400
  \item ￦500
  \item ￦800
  \item ￦1,000
\end{choices}
\end{question}
\answer{5}
